\documentclass[aspectratio=169,10pt]{beamer}

% Shared Beamer setup for Fall 2026 course slides.
% Clean Cal Poly-inspired theme: no custom class files or uncommon packages.
\mode<presentation>{
  \usetheme{default}
  \usefonttheme{professionalfonts}
}

\definecolor{CalPolyGreen}{HTML}{154734}
\definecolor{CalPolyGold}{HTML}{BD8B13}
\definecolor{CalPolySage}{HTML}{7FA28F}
\definecolor{CalPolyMint}{HTML}{DDF1E7}
\definecolor{CalPolySand}{HTML}{A29F86}
\definecolor{CalPolyBeige}{HTML}{EFEEDE}
\definecolor{DeckInk}{HTML}{1F2933}
\definecolor{DeckMuted}{HTML}{5F6B7A}
\definecolor{DeckSoft}{HTML}{F7F7F5}

\setbeamersize{text margin left=0.55in,text margin right=0.55in}
\setbeamercolor{normal text}{fg=DeckInk,bg=white}
\setbeamercolor{structure}{fg=CalPolyGreen}
\setbeamercolor{title}{fg=white}
\setbeamercolor{frametitle}{fg=CalPolyGreen,bg=white}
\setbeamercolor{block title}{bg=CalPolySage,fg=white}
\setbeamercolor{block body}{bg=CalPolyMint,fg=black}
\setbeamercolor{alerted text}{fg=CalPolyGold}

\setbeamerfont{title}{size=\Huge,series=\bfseries}
\setbeamerfont{frametitle}{size=\Large,series=\bfseries}
\setbeamerfont{block title}{size=\normalsize,series=\bfseries}

\setbeamertemplate{navigation symbols}{}
\setbeamertemplate{itemize item}{\color{CalPolyGold}$\blacktriangleright$}
\setbeamertemplate{itemize subitem}{\color{CalPolyGreen}$\bullet$}
\setbeamertemplate{footline}{
  \hfill{\color{DeckMuted}\scriptsize\insertframenumber/\inserttotalframenumber}
  \hspace{0.18in}\vspace{0.12in}
}
\setbeamertemplate{frametitle}{
  \vspace{0.18in}
  \hspace{0.02in}\insertframetitle\par
  \vspace{0.08in}
  \noindent\makebox[\linewidth][l]{%
    {\color{CalPolyGold}\rule{0.55in}{2pt}}%
    {\color{CalPolyGreen}\rule{\dimexpr\linewidth-0.55in\relax}{2pt}}%
  }
}

\newcommand{\CourseNumber}{Course}
\newcommand{\CourseTitle}{Course Title}
\newcommand{\LectureNumber}{0}
\newcommand{\LectureTitle}{Lecture Title}
\newcommand{\LectureDate}{Date}
\newcommand{\CourseUnit}{Unit}

\newcommand{\coursenumber}[1]{\renewcommand{\CourseNumber}{#1}}
\newcommand{\coursetitle}[1]{\renewcommand{\CourseTitle}{#1}}
\newcommand{\lecturenumber}[1]{\renewcommand{\LectureNumber}{#1}}
\newcommand{\lecturetitle}[1]{\renewcommand{\LectureTitle}{#1}}
\newcommand{\lecturedate}[1]{\renewcommand{\LectureDate}{#1}}
\newcommand{\courseunit}[1]{\renewcommand{\CourseUnit}{#1}}

\author{Kirk Duran}
\institute{California Polytechnic State University, San Luis Obispo}
\date{\LectureDate}

\newcommand{\makedecktitle}{
  \begingroup
  \setbeamercolor{background canvas}{bg=CalPolyGreen}
  \begin{frame}[plain]
    \vfill
    \begin{center}
      {\usebeamerfont{title}\color{white}\LectureTitle\par}
      \vspace{0.18in}
      {\Large\color{white}Lecture \LectureNumber\par}
    \end{center}
    \vfill
    \noindent{\color{white}\rule{\linewidth}{1.2pt}}\par
    \vspace{0.08in}
    \noindent\colorbox{CalPolyGold}{\makebox[0.24\linewidth][c]{\color{white}\Large\CourseNumber}}
    \hspace{0.12in}{\color{white}\Large Kirk Duran}
  \end{frame}
  \endgroup
}

\newenvironment{keyidea}{\begin{block}{Key idea}}{\end{block}}
\newenvironment{activity}{\begin{block}{In class}}{\end{block}}
\newenvironment{cpdefinition}{\begin{block}{Definition}}{\end{block}}
\newenvironment{exampleblockcp}{
  \begingroup
  \setbeamercolor{block title}{bg=CalPolySand,fg=white}
  \setbeamercolor{block body}{bg=CalPolyBeige,fg=black}
  \begin{block}{Example}
}{
  \end{block}
  \endgroup
}

\newcommand{\imageplaceholder}[2][0.3in]{%
  \noindent\makebox[\linewidth][c]{%
    \fcolorbox{CalPolySage}{CalPolyBeige}{%
      \parbox[c][#1][c]{0.86\linewidth}{%
        \centering
        {\color{DeckMuted}\scriptsize Image placeholder: }%
        {\color{DeckInk}\scriptsize\ttfamily\detokenize{#2}}%
      }%
    }%
  }%
}

\newcommand{\sectionframe}[1]{
  \begingroup
  \setbeamercolor{background canvas}{bg=CalPolyGreen}
  \begin{frame}[plain]
    \vfill
    {\color{white}\LARGE\bfseries #1\par}
    \vspace{0.18in}
    {\color{CalPolyGold}\rule{0.22\paperwidth}{3pt}}
    {\color{white}\rule{0.62\paperwidth}{3pt}}
    \vfill
  \end{frame}
  \endgroup
}

\newcommand{\placeholdernote}{
  \begin{block}{Placeholder}
    Replace these bullets with the final lecture material, examples, and in-class activities.
  \end{block}
}


\pdfstringdefDisableCommands{\def\translate#1{#1}}

\graphicspath{{../assets/lecture-20/}}

% If an image file is not yet available, the slide displays a labeled
% placeholder using the intended filename. Place the matching file in
% ../assets/lecture-20/ to replace the placeholder automatically.
\newcommand{\lectureimage}[2][1.75in]{%
  \IfFileExists{../assets/lecture-20/#2}{%
    \includegraphics[width=\linewidth,height=#1,keepaspectratio]{#2}%
  }{%
    \fbox{%
      \parbox[c][#1][c]{0.94\linewidth}{%
        \centering
        \small\textbf{Image Placeholder}\\[0.08in]
        \scriptsize\texttt{\detokenize{#2}}
      }%
    }%
  }%
}

\setbeamersize{text margin left=0.42in,text margin right=0.42in}

\coursenumber{CS 1001}
\coursetitle{Introduction to Computing}
\lecturenumber{20}
\lecturetitle{References, Aliasing, and Mutation}
\lecturedate{Fall 2026}
\courseunit{8. References and Mutation}

\begin{document}

\makedecktitle

\begin{frame}{Today}
  \begin{itemize}
    \item Review names, values, and immutable primitive data.
    \item Define mutability, identity, and aliasing precisely.
    \item Explain why assigning one list variable to another does not copy the list.
    \item Trace shared mutation, including mutation through function parameters.
    \item Contrast mutable lists with immutable strings and tuples.
    \item Create independent list copies when shared mutation is not desired.
  \end{itemize}

  \begin{keyidea}
    Assignment binds a name to an object. If two names refer to the same mutable object, a mutation through either name is visible through both.
  \end{keyidea}
\end{frame}

\sectionframe{Part 1: Values, Bindings, and Mutability}

\begin{frame}{Names Are Bound to Objects}
  \begin{block}{Example}
    \texttt{x = 10}\\
    \texttt{y = x}
  \end{block}

  \begin{itemize}
    \item Python evaluates the right-hand side, then binds the resulting object to the name on the left.
    \item A name is not the object itself.
    \item Multiple names may refer to the same object.
    \item This distinction becomes important when the shared object is mutable.
  \end{itemize}

  \begin{keyidea}
    When tracing references, track both the \textbf{name} and the \textbf{object} to which it refers.
  \end{keyidea}
\end{frame}

\begin{frame}{Mutable and Immutable Objects}
  \begin{cpdefinition}
    An object is \textbf{mutable} if its observable state can be changed after creation. An object is \textbf{immutable} if its value cannot be changed in place after creation.
  \end{cpdefinition}

  \begin{center}
    \renewcommand{\arraystretch}{1.3}
    \begin{tabular}{l|l}
      \textbf{Immutable in this unit} & \textbf{Mutable in this unit} \\\hline
      \texttt{int}, \texttt{float}, \texttt{bool} & \texttt{list} \\
      \texttt{str}, \texttt{tuple} &
    \end{tabular}
  \end{center}

  \begin{keyidea}
    Mutation changes an existing object. Reassignment changes which object a name refers to.
  \end{keyidea}
\end{frame}

\begin{frame}{Immutable Values: Rebinding Instead of Mutation}
  \begin{columns}[T,onlytextwidth]
    \begin{column}{0.56\textwidth}
      \begin{block}{Code}
        \texttt{x = 10}\\
        \texttt{y = x}\\
        \texttt{x = x + 1}
      \end{block}

      \begin{block}{Afterward}
        \texttt{x} $\longrightarrow$ \boxed{11}\\
        \texttt{y} $\longrightarrow$ \boxed{10}
      \end{block}

      The integer \texttt{10} was not changed into \texttt{11}; a new value was produced and \texttt{x} was rebound.
    \end{column}
    \begin{column}{0.37\textwidth}
      \lectureimage[2.55in]{immutable-rebinding.png}
    \end{column}
  \end{columns}
\end{frame}

\begin{frame}{Lists Introduce a New Possibility: In-Place Change}
  \begin{block}{Code}
    \texttt{scores = [90, 88, 76]}\\
    \texttt{scores[1] = 100}
  \end{block}

  \begin{itemize}
    \item The name \texttt{scores} still refers to the same list object.
    \item The state of that list object has changed.
    \item This is mutation, not rebinding.
  \end{itemize}

  \begin{keyidea}
    Once mutable data is shared, the distinction between changing a binding and changing an object becomes observable.
  \end{keyidea}
\end{frame}


\sectionframe{Part 2: Aliasing}

\begin{frame}{Definition: Aliasing}
  \begin{cpdefinition}
    \textbf{Aliasing} occurs when two or more names refer to the same object.
  \end{cpdefinition}

  \begin{block}{Example}
    \texttt{scores = [90, 88, 76]}\\
    \texttt{backup = scores}
  \end{block}

  \begin{alertblock}{Important}
    \texttt{backup = scores} does not duplicate the list. It creates another binding to the same list object.
  \end{alertblock}
\end{frame}

\begin{frame}{Mental Model: Two Names, One List}
  \begin{columns}[T,onlytextwidth]
    \begin{column}{0.48\textwidth}
      \begin{block}{Bindings}
        \texttt{scores} $\searrow$\\[0.08in]
        \hspace{0.75in}\boxed{[90, 88, 76]}\\[-0.02in]
        \texttt{backup} $\nearrow$
      \end{block}

      \begin{itemize}
        \item one list object;
        \item two references;
        \item mutation is shared.
      \end{itemize}
    \end{column}
    \begin{column}{0.45\textwidth}
      \lectureimage[2.65in]{alias-two-names-one-list.png}
    \end{column}
  \end{columns}
\end{frame}

\begin{frame}{Trace Shared Mutation}
  \begin{center}
    \renewcommand{\arraystretch}{1.28}
    \begin{tabular}{l|l|l}
      \textbf{Statement} & \textbf{Bindings} & \textbf{Object state} \\\hline
      \texttt{a = [1, 2, 3]} & \texttt{a -> L1} & \texttt{L1 = [1, 2, 3]} \\
      \texttt{b = a} & \texttt{a -> L1, b -> L1} & \texttt{L1 = [1, 2, 3]} \\
      \texttt{b[0] = 9} & \texttt{a -> L1, b -> L1} & \texttt{L1 = [9, 2, 3]}
    \end{tabular}
  \end{center}

  \begin{keyidea}
    Do not say ``\texttt{b} changed \texttt{a}.'' The indexed assignment changed the shared list object.
  \end{keyidea}
\end{frame}

\begin{frame}{Equality and Identity Answer Different Questions}
  \begin{block}{Code}
    \texttt{a = [1, 2]}\\
    \texttt{b = a}\\
    \texttt{c = [1, 2]}
  \end{block}

  \begin{center}
    \renewcommand{\arraystretch}{1.2}
    \begin{tabular}{l|c|l}
      \textbf{Expression} & \textbf{Result} & \textbf{Question} \\\hline
      \texttt{a == b} & \texttt{True} & same contents? \\
      \texttt{a is b} & \texttt{True} & same object? \\
      \texttt{a == c} & \texttt{True} & same contents? \\
      \texttt{a is c} & \texttt{False} & same object?
    \end{tabular}
  \end{center}
\end{frame}

\begin{frame}{Prediction: Follow the Object, Not Just the Variable Name}
  \begin{block}{Code}
    \texttt{first = [4, 5, 6]}\\
    \texttt{second = first}\\
    \texttt{first[-1] = 0}
  \end{block}

  \begin{activity}
    Draw the reference diagram after each statement. Then determine the values observed through \texttt{first} and \texttt{second}.
  \end{activity}
\end{frame}


\sectionframe{Part 3: Mutation as an Observable Side Effect}

\begin{frame}{Function Parameters Can Alias Caller-Owned Lists}
  \begin{columns}[T,onlytextwidth]
    \begin{column}{0.57\textwidth}
      \begin{block}{Code}
        \texttt{def set\_first(values: list, new\_value: int) -> None:}\\
        \hspace{0.18in}\texttt{values[0] = new\_value}\\[0.08in]
        \texttt{scores = [70, 80, 90]}\\
        \texttt{set\_first(scores, 100)}
      \end{block}

      During the call, \texttt{scores} and parameter \texttt{values} refer to the same list.
    \end{column}
    \begin{column}{0.36\textwidth}
      \lectureimage[2.55in]{parameter-alias-call-frame.png}
    \end{column}
  \end{columns}
\end{frame}

\begin{frame}{Definition: Side Effect}
  \begin{cpdefinition}
    A \textbf{side effect} is an observable change to program state that occurs in addition to producing a return value. Here, our main example is mutation of a list supplied by the caller.
  \end{cpdefinition}

  \begin{itemize}
    \item \texttt{set\_first(...)} returns \texttt{None}, but changes the caller-visible list.
    \item Methods such as \texttt{append()}, \texttt{remove()}, and \texttt{reverse()} also intentionally mutate lists.
    \item Mutation is not inherently an error; unexpected or undocumented shared mutation is the problem.
  \end{itemize}
\end{frame}

\begin{frame}{Trace a Function Side Effect}
  \begin{block}{Code}
    \texttt{def replace\_last(items: list, value: int) -> None:}\\
    \hspace{0.18in}\texttt{items[-1] = value}\\[0.08in]
    \texttt{data = [2, 4, 6]}\\
    \texttt{replace\_last(data, 99)}
  \end{block}

  \begin{activity}
    At the indexed assignment, identify every active name referring to the list object. What is the final value observed through \texttt{data}?
  \end{activity}
\end{frame}


\sectionframe{Part 4: Immutable Sequences — Strings and Tuples}

\begin{frame}{Strings Prevent In-Place Element Mutation}
  \begin{block}{Code}
    \texttt{text = "cat"}\\
    \texttt{text[0] = "b"}
  \end{block}

  \begin{alertblock}{Result}
    Python raises \texttt{TypeError}. A string does not support indexed assignment.
  \end{alertblock}

  \begin{itemize}
    \item Indexing can read a string position.
    \item Slicing and concatenation create new strings.
    \item The existing string value cannot be changed in place.
  \end{itemize}
\end{frame}

\begin{frame}{String ``Change'' Means Construct, Then Rebind}
  \begin{columns}[T,onlytextwidth]
    \begin{column}{0.56\textwidth}
      \begin{block}{Code}
        \texttt{a = "cat"}\\
        \texttt{b = a}\\
        \texttt{a = "b" + a[1:]}
      \end{block}

      \begin{block}{Afterward}
        \texttt{a} $\longrightarrow$ \texttt{"bat"}\\
        \texttt{b} $\longrightarrow$ \texttt{"cat"}
      \end{block}

      The original \texttt{"cat"} object was not mutated.
    \end{column}
    \begin{column}{0.37\textwidth}
      \lectureimage[2.55in]{immutable-string-rebinding.png}
    \end{column}
  \end{columns}
\end{frame}

\begin{frame}{Definition: Tuple}
  \begin{cpdefinition}
    A \textbf{tuple} is an ordered, indexed, immutable sequence of references.
  \end{cpdefinition}

  \begin{block}{Examples}
    \texttt{point = (3, 5)}\qquad
    \texttt{empty = ()}\qquad
    \texttt{single = (42,)}
  \end{block}

  \begin{center}
    \renewcommand{\arraystretch}{1.2}
    \begin{tabular}{l|l}
      \texttt{point[0]} & \texttt{3} \\
      \texttt{point[:1]} & \texttt{(3,)} \\
      \texttt{len(point)} & \texttt{2} \\
      \texttt{5 in point} & \texttt{True}
    \end{tabular}
  \end{center}
\end{frame}

\begin{frame}{To ``Change'' a Tuple, Build a New Tuple}
  \begin{block}{Invalid mutation}
    \texttt{data = (10, 20, 30)}\\
    \texttt{data[1] = 99}\qquad $\rightarrow$ \quad \texttt{TypeError}
  \end{block}

  \begin{block}{Construct and rebind}
    \texttt{data = data[:1] + (99,) + data[2:]}\\
    \texttt{data} $\longrightarrow$ \texttt{(10, 99, 30)}
  \end{block}

  \begin{keyidea}
    Slicing and concatenation produce a new tuple; assignment then rebinds the name.
  \end{keyidea}
\end{frame}

\begin{frame}{Lists, Strings, and Tuples Share a Sequence Model}
  \begin{center}
    \renewcommand{\arraystretch}{1.3}
    \begin{tabular}{l|c|c|c}
      & \textbf{List} & \textbf{String} & \textbf{Tuple} \\\hline
      Ordered / indexed & yes & yes & yes \\
      Sliceable & yes & yes & yes \\
      \texttt{len()} / membership & yes & yes & yes \\
      Element assignment & yes & no & no \\
      Mutable & yes & no & no
    \end{tabular}
  \end{center}
\end{frame}


\sectionframe{Part 5: Creating Independent List Copies}

\begin{frame}{When Aliasing Is Not What You Want}
  \begin{block}{Problem}
    \texttt{original = [10, 20, 30]}\\
    \texttt{working = original}\\
    \texttt{working[0] = 99}
  \end{block}

  \begin{itemize}
    \item \texttt{working} is an alias, not an independent working copy.
    \item Mutation is therefore visible through \texttt{original}.
    \item If independent mutation is required, deliberately create a new list object.
  \end{itemize}
\end{frame}

\begin{frame}{\texttt{copy()} Creates a New List Object}
  \begin{columns}[T,onlytextwidth]
    \begin{column}{0.55\textwidth}
      \begin{block}{Code}
        \texttt{original = [10, 20, 30]}\\
        \texttt{working = original.copy()}\\
        \texttt{working[0] = 99}
      \end{block}

      \begin{block}{Afterward}
        \texttt{original} $\longrightarrow$ \texttt{[10, 20, 30]}\\
        \texttt{working} $\longrightarrow$ \texttt{[99, 20, 30]}
      \end{block}
    \end{column}
    \begin{column}{0.38\textwidth}
      \lectureimage[2.60in]{alias-versus-copy.png}
    \end{column}
  \end{columns}
\end{frame}

\begin{frame}{Other Ways to Create an Independent List Container}
  \begin{block}{Full slice}
    \texttt{working = original[:]}
  \end{block}

  \begin{block}{List construction}
    \texttt{working = list(original)}
  \end{block}

  \begin{itemize}
    \item Both forms create a new list object.
    \item The new list initially contains equal element values.
    \item For the primitive-value lists used in this unit, mutation of positions in one list does not alter the other list.
  \end{itemize}
\end{frame}

\begin{frame}{Alias vs. Copy: Equality Can Be the Same While Identity Differs}
  \begin{block}{Code}
    \texttt{a = [1, 2, 3]}\\
    \texttt{b = a}\\
    \texttt{c = a.copy()}
  \end{block}

  \begin{center}
    \renewcommand{\arraystretch}{1.2}
    \begin{tabular}{l|c@{\qquad}l|c}
      \texttt{a == b} & \texttt{True} & \texttt{a is b} & \texttt{True} \\
      \texttt{a == c} & \texttt{True} & \texttt{a is c} & \texttt{False}
    \end{tabular}
  \end{center}

  \begin{keyidea}
    Copying can preserve current contents while intentionally changing object identity.
  \end{keyidea}
\end{frame}

\begin{frame}{Copy Before Mutation When Isolation Is Required}
  \begin{block}{Example}
    \texttt{def changed\_first(values: list, new\_value: int) -> list:}\\
    \hspace{0.18in}\texttt{result = values.copy()}\\
    \hspace{0.18in}\texttt{result[0] = new\_value}\\
    \hspace{0.18in}\texttt{return result}
  \end{block}

  \begin{itemize}
    \item The parameter still refers to the caller's list.
    \item Mutation is performed only on the new list.
    \item The caller's original list remains unchanged.
  \end{itemize}
\end{frame}

\begin{frame}{Two Valid Function Designs}
  \begin{columns}[T,onlytextwidth]
    \begin{column}{0.47\textwidth}
      \begin{block}{Mutate caller-owned list}
        \texttt{def set\_first(xs, value):}\\
        \hspace{0.14in}\texttt{xs[0] = value}
      \end{block}
      \begin{itemize}
        \item intentional side effect;
        \item returns \texttt{None}.
      \end{itemize}
    \end{column}
    \begin{column}{0.47\textwidth}
      \begin{block}{Return a modified copy}
        \texttt{def with\_first(xs, value):}\\
        \hspace{0.14in}\texttt{result = xs.copy()}\\
        \hspace{0.14in}\texttt{result[0] = value}\\
        \hspace{0.14in}\texttt{return result}
      \end{block}
      \begin{itemize}
        \item original stays unchanged;
        \item new list is returned.
      \end{itemize}
    \end{column}
  \end{columns}
\end{frame}


\sectionframe{Part 6: Integrated Tracing and Design Rules}

\begin{frame}{Integrated Trace: Alias and Independent Copy}
  \begin{block}{Code}
    \texttt{a = [10, 20, 30]}\\
    \texttt{b = a}\\
    \texttt{c = a.copy()}\\
    \texttt{b[1] = 99}\\
    \texttt{c[0] = 5}
  \end{block}

  \begin{activity}
    Determine the final values observed through \texttt{a}, \texttt{b}, and \texttt{c}. Draw the reference diagram after each statement.
  \end{activity}
\end{frame}

\begin{frame}{Integrated Trace: Immutable Rebinding}
  \begin{block}{Code}
    \texttt{s1 = "code"}\\
    \texttt{s2 = s1}\\
    \texttt{s1 = "C" + s1[1:]}\\[0.08in]
    \texttt{t1 = (1, 2, 3)}\\
    \texttt{t2 = t1}\\
    \texttt{t1 = t1[:1] + (9,) + t1[2:]}
  \end{block}

  \begin{activity}
    Determine the final values of all four names. Identify every newly constructed object.
  \end{activity}
\end{frame}

\begin{frame}{Common Incorrect Mental Models}
  \begin{center}
    \renewcommand{\arraystretch}{1.3}
    \begin{tabular}{p{0.43\textwidth}|p{0.46\textwidth}}
      \textbf{Incorrect} & \textbf{Replace it with} \\\hline
      ``Each variable stores its own copy.'' & A name refers to an object; names may alias one object. \\
      ``Changing \texttt{b} changed \texttt{a}.'' & Mutation changed the shared object. \\
      ``Reassignment mutates the old value.'' & Reassignment changes a binding. \\
      ``A string method edits the string.'' & String operations produce new string values.
    \end{tabular}
  \end{center}
\end{frame}

\begin{frame}{A Practical Tracing Procedure}
  \begin{enumerate}
    \item Identify each object that is created.
    \item Record which names refer to each object.
    \item For assignment, decide whether the result is a new object or an existing reference.
    \item For a mutating operation, update the object itself.
    \item For an immutable operation, record the newly produced value and any rebinding.
    \item During a function call, check whether a parameter aliases caller-owned mutable data.
  \end{enumerate}
\end{frame}

\begin{frame}{Choosing Between a List and a Tuple}
  \begin{center}
    \renewcommand{\arraystretch}{1.35}
    \begin{tabular}{p{0.43\textwidth}|p{0.43\textwidth}}
      \textbf{Use a list when...} & \textbf{Use a tuple when...} \\\hline
      positions are expected to change & the sequence should remain fixed \\
      elements may be inserted or removed & element replacement should be disallowed \\
      mutation is part of the intended model & immutability communicates the intended model
    \end{tabular}
  \end{center}
\end{frame}

\begin{frame}{Operation Summary}
  \begin{center}
    \renewcommand{\arraystretch}{1.22}
    \begin{tabular}{l|l}
      \textbf{Operation} & \textbf{Effect} \\\hline
      \texttt{b = a} where \texttt{a} is a list & alias existing list \\
      \texttt{a[i] = value} & mutate list object \\
      \texttt{a.copy()}, \texttt{a[:]}, \texttt{list(a)} & create new list \\
      string indexed assignment & invalid \\
      tuple indexed assignment & invalid \\
      string/tuple reconstruction + assignment & create new value, then rebind
    \end{tabular}
  \end{center}
\end{frame}

\begin{frame}{Takeaways}
  \begin{itemize}
    \item Names refer to objects; assignment does not automatically copy mutable data.
    \item Aliasing means multiple names refer to the same object.
    \item Mutation changes a mutable object, so every alias observes the new state.
    \item Function parameters can alias caller-owned lists, making mutation visible outside the call.
    \item Integers, floats, booleans, strings, and tuples are immutable; apparent changes require new values and rebinding.
    \item Tuples retain sequence operations while preventing element replacement.
    \item Use \texttt{copy()}, a full slice, or \texttt{list(...)} when an independent list is required.
  \end{itemize}

  \begin{keyidea}
    When tracing mutable data, ask not only ``What value does this variable show?'' but also ``Which object does this name refer to?''
  \end{keyidea}
\end{frame}

\end{document}
